% ==============================================================================
% CONFIG.TEX - CẤU HÌNH ĐỊNH DẠNG BÁO CÁO THEO QUY CHUẨN ICTU
% ==============================================================================

\usepackage[utf8]{inputenc}
\usepackage[vietnamese]{babel}

% 1. Cấu hình Font chữ Times New Roman
\usepackage{iftex}
\ifPDFTeX
    \usepackage[T5,T1]{fontenc}
    \usepackage{mathptmx} % Times New Roman cho pdfLaTeX
\else
    \usepackage{fontspec}
    \setmainfont{Times New Roman} % Times New Roman hệ thống cho XeLaTeX / LuaLaTeX
\fi

% 2. Cỡ chữ 13pt (Quy định: Times New Roman cỡ chữ 13)
\usepackage{scrextend}
\changefontsizes[17pt]{13pt}

% 3. Căn lề khổ giấy A4 (Quy định: Trên 2.5cm, Dưới 2.5cm, Trái 3.0cm, Phải 2.0cm)
\usepackage[
    a4paper,
    top=2.5cm,
    bottom=2.5cm,
    left=3.0cm,
    right=2.0cm,
    footskip=1.2cm
]{geometry}

% 4. Dãn dòng và Dãn đoạn (Quy định: Multiple 1.3 lines; Before 6pt, After 6pt)
\usepackage{setspace}
\setstretch{1.3} % Multiple 1.3 lines
\setlength{\parskip}{6pt} % Before/After 6pt
\setlength{\parindent}{1.0cm} % Thụt đầu dòng 1cm

% 5. Đánh số trang (Quy định: Đánh ở giữa, phía dưới mỗi trang giấy)
\usepackage{fancyhdr}
\pagestyle{fancy}
\fancyhf{}
\cfoot{\thepage}
\renewcommand{\headrulewidth}{0pt}
\renewcommand{\footrulewidth}{0pt}

% Cấu hình trang plain (đầu chương) vẫn giữ số trang ở giữa dưới
\fancypagestyle{plain}{
    \fancyhf{}
    \cfoot{\thepage}
    \renewcommand{\headrulewidth}{0pt}
    \renewcommand{\footrulewidth}{0pt}
}

% 6. Đánh số tiểu mục (Quy định: Tối đa 4 chữ số, X.Y.Z.W)
\setcounter{secnumdepth}{4}
\setcounter{tocdepth}{3}

% 7. Định dạng tiêu đề chương và các mục (Quy định: Tên các chương VIẾT HOA, ĐẬM)
\usepackage{titlesec}

% Định dạng Chương
\titleformat{\chapter}[display]
  {\normalfont\bfseries\centering}
  {\MakeUppercase{\chaptertitlename\ \thechapter}}
  {8pt}
  {\LARGE\bfseries\MakeUppercase}
\titlespacing*{\chapter}{0pt}{-10pt}{18pt}

% Định dạng Section (Cấp 1: 1.1)
\titleformat{\section}
  {\normalfont\fontsize{14pt}{18pt}\bfseries}
  {\thesection}{1em}{}
\titlespacing*{\section}{0pt}{12pt}{6pt}

% Định dạng Subsection (Cấp 2: 1.1.1)
\titleformat{\subsection}
  {\normalfont\fontsize{13pt}{17pt}\bfseries}
  {\thesubsection}{1em}{}
\titlespacing*{\subsection}{0pt}{10pt}{4pt}

% Định dạng Subsubsection (Cấp 3: 1.1.1.1)
\titleformat{\subsubsection}
  {\normalfont\fontsize{13pt}{17pt}\bfseries\itshape}
  {\thesubsubsection}{1em}{}
\titlespacing*{\subsubsection}{0pt}{8pt}{3pt}

% 8. Định dạng Bảng biểu và Hình ảnh (Quy định: Đầu đề Bảng ở TRÊN, đầu đề Hình ở DƯỚI)
\usepackage{graphicx}
\graphicspath{{figures/}}
\newcommand{\figpath}{figures/}
\usepackage{caption}
\usepackage{subcaption}
\captionsetup{
    font={small},
    labelfont=bf,
    labelsep=colon
}
\captionsetup[table]{
    position=top,
    skip=6pt
}
\captionsetup[figure]{
    position=bottom,
    skip=6pt
}

% 9. Các gói bổ trợ toán học, bảng biểu, liên kết và thuật toán
\usepackage{amsmath,amssymb,amsfonts}
\usepackage{booktabs}
\usepackage{array}
\usepackage{tabularx}
\usepackage{multirow}
\usepackage{listings}
\usepackage{xcolor}
\usepackage{tikz}
\usetikzlibrary{calc}
\usepackage{hyperref}

\hypersetup{
    colorlinks=true,
    linkcolor=black,
    citecolor=black,
    urlcolor=blue,
    plainpages=false,
    pdfpagelabels,
    pdfauthor={Nhóm 2 - Điện toán đám mây},
    pdftitle={Báo cáo Bài tập lớn: Quy trình CI/CD và Triển khai hệ thống nhận diện ngôn ngữ ký hiệu tiếng Việt}
}

% Cấu hình hiển thị code
\lstset{
    basicstyle=\ttfamily\small,
    keywordstyle=\color{blue}\bfseries,
    stringstyle=\color{red},
    commentstyle=\color{gray}\itshape,
    numbers=left,
    numberstyle=\tiny\color{gray},
    stepnumber=1,
    numbersep=8pt,
    backgroundcolor=\color{white},
    showspaces=false,
    showstringspaces=false,
    showtabs=false,
    frame=single,
    tabsize=2,
    captionpos=b,
    breaklines=true,
    breakatwhitespace=false
}
